\documentclass[11pt,a4paper]{article}
\newcommand{\gradelabel}{Grade 10}
\newcommand{\leconhdr}{Conditions}
% Préambule commun aux plans de cours de robotique (Grades 7 à 12).
% Chaque plan de cours commence par \documentclass[11pt,a4paper]{article}, définit
% \gradelabel et \leconhdr (en-têtes de page), puis % Préambule commun aux plans de cours de robotique (Grades 7 à 12).
% Chaque plan de cours commence par \documentclass[11pt,a4paper]{article}, définit
% \gradelabel et \leconhdr (en-têtes de page), puis % Préambule commun aux plans de cours de robotique (Grades 7 à 12).
% Chaque plan de cours commence par \documentclass[11pt,a4paper]{article}, définit
% \gradelabel et \leconhdr (en-têtes de page), puis \input{../../preambule_robotique.tex}.
\usepackage[T1]{fontenc}
\usepackage[utf8]{inputenc}
\usepackage{lmodern}
\usepackage[french]{babel}
\usepackage[margin=2cm,headheight=14pt]{geometry}
\usepackage[most]{tcolorbox}
\usepackage{xcolor,tabularx,booktabs,array,enumitem,fancyhdr,amssymb,tikz}
\usetikzlibrary{arrows.meta,positioning}

\pagestyle{fancy}
\fancyhf{}
\lhead{\small Robotique -- \gradelabel{} -- Novembre 2026}
\rhead{\small \leconhdr}
\cfoot{\thepage}
\setlist{itemsep=1pt,topsep=2pt,leftmargin=*}
\renewcommand{\arraystretch}{1.35}
\newcolumntype{Y}{>{\raggedright\arraybackslash}X}
\newcommand{\cb}{$\square$\,}
\newcommand{\fl}{$\rightarrow$}
\newcommand{\lig}{\noindent\rule{\linewidth}{0.3pt}\\[5mm]}

\newtcolorbox{etape}[3]{enhanced,breakable,colback=#1!6,colframe=#1!70!black,
  coltitle=white,fonttitle=\bfseries,title={#2\hfill #3},boxrule=0.8pt,arc=2mm,
  before skip=8pt,after skip=8pt}
\newtcolorbox{astuce}[1][Différenciation et besoins particuliers]{breakable,
  colback=yellow!12,colframe=yellow!50!black,fonttitle=\bfseries\small,
  title={#1},boxrule=0.5pt,arc=1mm,before skip=6pt,after skip=6pt}
\newtcolorbox{code}[1][Code]{breakable,colback=black!4,colframe=black!60,
  fonttitle=\bfseries\small,title={#1},boxrule=0.5pt,arc=1mm,fontupper=\ttfamily\small,
  before skip=6pt,after skip=6pt}
% Blocs de code coloré (écrire le code sans accents : comme dans l'éditeur)
\newtcblisting{arduino}[1][Code Arduino]{listing only,breakable,colback=black!4,
  colframe=black!60,fonttitle=\bfseries\small,title={#1},boxrule=0.5pt,arc=1mm,
  before skip=6pt,after skip=6pt,
  listing options={language=C++,basicstyle=\ttfamily\small,columns=fullflexible,
  keepspaces=true,showstringspaces=false,commentstyle=\color{green!40!black}\itshape,
  keywordstyle=\color{blue!70!black}\bfseries}}
\newtcblisting{python}[1][Code Python]{listing only,breakable,colback=black!4,
  colframe=black!60,fonttitle=\bfseries\small,title={#1},boxrule=0.5pt,arc=1mm,
  before skip=6pt,after skip=6pt,
  listing options={language=Python,basicstyle=\ttfamily\small,columns=fullflexible,
  keepspaces=true,showstringspaces=false,commentstyle=\color{green!40!black}\itshape,
  keywordstyle=\color{blue!70!black}\bfseries}}
\newcommand{\entete}[2]{%
  \begin{tcolorbox}[colback=#1!10,colframe=#1!70!black,boxrule=0.8pt,arc=2mm]
  {\Large\bfseries #2}\hfill Nom : \rule{4cm}{0.3pt}\quad Date : \rule{2.2cm}{0.3pt}
  \end{tcolorbox}}

% Bandeau de titre : #1 titre, #2 semaine
\newcommand{\titreblock}[2]{%
\begin{tcolorbox}[colback=teal!10,colframe=teal!70!black,arc=2mm,boxrule=1pt]
\centering
{\LARGE\bfseries #1}\\[2mm]
{\large Robotique -- \gradelabel{} -- Plan de cours de 2 heures}\\[1mm]
Enseignant : Francis Ki \quad|\quad #2
\end{tcolorbox}}

% Tableau du minutage : 6 contenus (accroche, mini-leçon, guidée, autonome, mise en commun, exit ticket)
\newcommand{\minutage}[6]{%
\begin{tabularx}{\linewidth}{>{\bfseries}p{2.2cm}p{2.6cm}Y p{3cm}}
\toprule
Minutes & Étape & Contenu & Organisation\\
\midrule
0 -- 10 & Accroche & #1 & Classe entière\\
10 -- 30 & Mini-leçon & #2 & Classe entière\\
30 -- 60 & Pratique guidée & #3 & Collectif puis binômes\\
60 -- 95 & Pratique autonome & #4 & Individuel sur poste\\
95 -- 110 & Mise en commun & #5 & Classe entière\\
110 -- 120 & Exit ticket & #6 & Individuel sur poste\\
\bottomrule
\end{tabularx}}

% Étayage progressif : #1 aide 1, #2 aide 2, #3 aide 3
\newcommand{\aides}[3]{%
\textbf{Étayage progressif de l'aide.}
\begin{center}
\begin{tabularx}{\linewidth}{>{\bfseries}p{2.4cm}Y p{4.4cm}}
\toprule
Niveau d'aide & Contenu & Qui intervient\\
\midrule
Aide 1 & #1 & Le binôme seul\\
Aide 2 & #2 & Le binôme avec la carte d'aide\\
Aide 3 & #3 & L'enseignant\\
\bottomrule
\end{tabularx}
\end{center}}

% Les trois groupes de la pratique autonome : #1 tâche A, #2 tâche B, #3 tâche C
\newcommand{\groupes}[3]{%
\textbf{Consigne commune\,:} chaque élève ouvre le fichier de son groupe, le renomme avec son prénom et travaille seul sur son poste. L'enseignant passe d'abord auprès du groupe A pendant 10 min.

\medskip
\begin{tabularx}{\linewidth}{>{\bfseries}p{3cm}Y p{2.2cm}}
\toprule
Groupe & Tâche sur poste & Annexe\\
\midrule
A : en dessous du niveau & #1 & Annexe C\\
B : au niveau & #2 & Annexe D\\
C : au-dessus du niveau & #3 & Annexe E\\
\bottomrule
\end{tabularx}}

% Mise en commun : #1 groupe A, #2 groupe B, #3 groupe C, #4 conclusion
\newcommand{\miseencommun}[4]{%
\begin{etape}{teal}{5. Mise en commun}{95 -- 110 min}
\begin{enumerate}
 \item (2 min) Retour au calme, écrans mis en veille.
 \item (9 min) Trois passages courts de 3 min, avec valorisation :
 \begin{itemize}
  \item \textbf{Groupe A} : #1
  \item \textbf{Groupe B} : #2
  \item \textbf{Groupe C} : #3
 \end{itemize}
 \item (4 min) Conclusion : #4
\end{enumerate}
\end{etape}}

% Clôture : #1 question 1, #2 question 2
\newcommand{\cloture}[2]{%
\begin{etape}{gray}{6. Clôture : exit ticket}{110 -- 120 min}
Formulaire en ligne ou papier (annexe G) :
\begin{enumerate}
 \item #1
 \item #2
 \item Auto-évaluation : \emph{Je maîtrise / Je progresse / J'ai besoin d'aide.}
\end{enumerate}
Lire les réponses avant la séance suivante pour ajuster les groupes de besoin.
\end{etape}}

% Grille d'observation : #1 #2 #3 intitulés de colonnes
\newcommand{\grille}[3]{%
\subsection*{Grille d'observation (enseignant)}
\begin{tabularx}{\linewidth}{Y p{2.4cm} p{2.4cm} p{2.4cm} Y}
\toprule
Élève & #1 & #2 & #3 & Groupe suivant\\
\midrule
 & & & & \\ & & & & \\ & & & & \\ & & & & \\ & & & & \\
\bottomrule
\end{tabularx}}

% Signature de l'exit ticket
\newcommand{\autoeval}{Je me sens : \cb Je maîtrise \quad \cb Je progresse \quad \cb J'ai besoin d'aide}
.
\usepackage[T1]{fontenc}
\usepackage[utf8]{inputenc}
\usepackage{lmodern}
\usepackage[french]{babel}
\usepackage[margin=2cm,headheight=14pt]{geometry}
\usepackage[most]{tcolorbox}
\usepackage{xcolor,tabularx,booktabs,array,enumitem,fancyhdr,amssymb,tikz}
\usetikzlibrary{arrows.meta,positioning}

\pagestyle{fancy}
\fancyhf{}
\lhead{\small Robotique -- \gradelabel{} -- Novembre 2026}
\rhead{\small \leconhdr}
\cfoot{\thepage}
\setlist{itemsep=1pt,topsep=2pt,leftmargin=*}
\renewcommand{\arraystretch}{1.35}
\newcolumntype{Y}{>{\raggedright\arraybackslash}X}
\newcommand{\cb}{$\square$\,}
\newcommand{\fl}{$\rightarrow$}
\newcommand{\lig}{\noindent\rule{\linewidth}{0.3pt}\\[5mm]}

\newtcolorbox{etape}[3]{enhanced,breakable,colback=#1!6,colframe=#1!70!black,
  coltitle=white,fonttitle=\bfseries,title={#2\hfill #3},boxrule=0.8pt,arc=2mm,
  before skip=8pt,after skip=8pt}
\newtcolorbox{astuce}[1][Différenciation et besoins particuliers]{breakable,
  colback=yellow!12,colframe=yellow!50!black,fonttitle=\bfseries\small,
  title={#1},boxrule=0.5pt,arc=1mm,before skip=6pt,after skip=6pt}
\newtcolorbox{code}[1][Code]{breakable,colback=black!4,colframe=black!60,
  fonttitle=\bfseries\small,title={#1},boxrule=0.5pt,arc=1mm,fontupper=\ttfamily\small,
  before skip=6pt,after skip=6pt}
% Blocs de code coloré (écrire le code sans accents : comme dans l'éditeur)
\newtcblisting{arduino}[1][Code Arduino]{listing only,breakable,colback=black!4,
  colframe=black!60,fonttitle=\bfseries\small,title={#1},boxrule=0.5pt,arc=1mm,
  before skip=6pt,after skip=6pt,
  listing options={language=C++,basicstyle=\ttfamily\small,columns=fullflexible,
  keepspaces=true,showstringspaces=false,commentstyle=\color{green!40!black}\itshape,
  keywordstyle=\color{blue!70!black}\bfseries}}
\newtcblisting{python}[1][Code Python]{listing only,breakable,colback=black!4,
  colframe=black!60,fonttitle=\bfseries\small,title={#1},boxrule=0.5pt,arc=1mm,
  before skip=6pt,after skip=6pt,
  listing options={language=Python,basicstyle=\ttfamily\small,columns=fullflexible,
  keepspaces=true,showstringspaces=false,commentstyle=\color{green!40!black}\itshape,
  keywordstyle=\color{blue!70!black}\bfseries}}
\newcommand{\entete}[2]{%
  \begin{tcolorbox}[colback=#1!10,colframe=#1!70!black,boxrule=0.8pt,arc=2mm]
  {\Large\bfseries #2}\hfill Nom : \rule{4cm}{0.3pt}\quad Date : \rule{2.2cm}{0.3pt}
  \end{tcolorbox}}

% Bandeau de titre : #1 titre, #2 semaine
\newcommand{\titreblock}[2]{%
\begin{tcolorbox}[colback=teal!10,colframe=teal!70!black,arc=2mm,boxrule=1pt]
\centering
{\LARGE\bfseries #1}\\[2mm]
{\large Robotique -- \gradelabel{} -- Plan de cours de 2 heures}\\[1mm]
Enseignant : Francis Ki \quad|\quad #2
\end{tcolorbox}}

% Tableau du minutage : 6 contenus (accroche, mini-leçon, guidée, autonome, mise en commun, exit ticket)
\newcommand{\minutage}[6]{%
\begin{tabularx}{\linewidth}{>{\bfseries}p{2.2cm}p{2.6cm}Y p{3cm}}
\toprule
Minutes & Étape & Contenu & Organisation\\
\midrule
0 -- 10 & Accroche & #1 & Classe entière\\
10 -- 30 & Mini-leçon & #2 & Classe entière\\
30 -- 60 & Pratique guidée & #3 & Collectif puis binômes\\
60 -- 95 & Pratique autonome & #4 & Individuel sur poste\\
95 -- 110 & Mise en commun & #5 & Classe entière\\
110 -- 120 & Exit ticket & #6 & Individuel sur poste\\
\bottomrule
\end{tabularx}}

% Étayage progressif : #1 aide 1, #2 aide 2, #3 aide 3
\newcommand{\aides}[3]{%
\textbf{Étayage progressif de l'aide.}
\begin{center}
\begin{tabularx}{\linewidth}{>{\bfseries}p{2.4cm}Y p{4.4cm}}
\toprule
Niveau d'aide & Contenu & Qui intervient\\
\midrule
Aide 1 & #1 & Le binôme seul\\
Aide 2 & #2 & Le binôme avec la carte d'aide\\
Aide 3 & #3 & L'enseignant\\
\bottomrule
\end{tabularx}
\end{center}}

% Les trois groupes de la pratique autonome : #1 tâche A, #2 tâche B, #3 tâche C
\newcommand{\groupes}[3]{%
\textbf{Consigne commune\,:} chaque élève ouvre le fichier de son groupe, le renomme avec son prénom et travaille seul sur son poste. L'enseignant passe d'abord auprès du groupe A pendant 10 min.

\medskip
\begin{tabularx}{\linewidth}{>{\bfseries}p{3cm}Y p{2.2cm}}
\toprule
Groupe & Tâche sur poste & Annexe\\
\midrule
A : en dessous du niveau & #1 & Annexe C\\
B : au niveau & #2 & Annexe D\\
C : au-dessus du niveau & #3 & Annexe E\\
\bottomrule
\end{tabularx}}

% Mise en commun : #1 groupe A, #2 groupe B, #3 groupe C, #4 conclusion
\newcommand{\miseencommun}[4]{%
\begin{etape}{teal}{5. Mise en commun}{95 -- 110 min}
\begin{enumerate}
 \item (2 min) Retour au calme, écrans mis en veille.
 \item (9 min) Trois passages courts de 3 min, avec valorisation :
 \begin{itemize}
  \item \textbf{Groupe A} : #1
  \item \textbf{Groupe B} : #2
  \item \textbf{Groupe C} : #3
 \end{itemize}
 \item (4 min) Conclusion : #4
\end{enumerate}
\end{etape}}

% Clôture : #1 question 1, #2 question 2
\newcommand{\cloture}[2]{%
\begin{etape}{gray}{6. Clôture : exit ticket}{110 -- 120 min}
Formulaire en ligne ou papier (annexe G) :
\begin{enumerate}
 \item #1
 \item #2
 \item Auto-évaluation : \emph{Je maîtrise / Je progresse / J'ai besoin d'aide.}
\end{enumerate}
Lire les réponses avant la séance suivante pour ajuster les groupes de besoin.
\end{etape}}

% Grille d'observation : #1 #2 #3 intitulés de colonnes
\newcommand{\grille}[3]{%
\subsection*{Grille d'observation (enseignant)}
\begin{tabularx}{\linewidth}{Y p{2.4cm} p{2.4cm} p{2.4cm} Y}
\toprule
Élève & #1 & #2 & #3 & Groupe suivant\\
\midrule
 & & & & \\ & & & & \\ & & & & \\ & & & & \\ & & & & \\
\bottomrule
\end{tabularx}}

% Signature de l'exit ticket
\newcommand{\autoeval}{Je me sens : \cb Je maîtrise \quad \cb Je progresse \quad \cb J'ai besoin d'aide}
.
\usepackage[T1]{fontenc}
\usepackage[utf8]{inputenc}
\usepackage{lmodern}
\usepackage[french]{babel}
\usepackage[margin=2cm,headheight=14pt]{geometry}
\usepackage[most]{tcolorbox}
\usepackage{xcolor,tabularx,booktabs,array,enumitem,fancyhdr,amssymb,tikz}
\usetikzlibrary{arrows.meta,positioning}

\pagestyle{fancy}
\fancyhf{}
\lhead{\small Robotique -- \gradelabel{} -- Novembre 2026}
\rhead{\small \leconhdr}
\cfoot{\thepage}
\setlist{itemsep=1pt,topsep=2pt,leftmargin=*}
\renewcommand{\arraystretch}{1.35}
\newcolumntype{Y}{>{\raggedright\arraybackslash}X}
\newcommand{\cb}{$\square$\,}
\newcommand{\fl}{$\rightarrow$}
\newcommand{\lig}{\noindent\rule{\linewidth}{0.3pt}\\[5mm]}

\newtcolorbox{etape}[3]{enhanced,breakable,colback=#1!6,colframe=#1!70!black,
  coltitle=white,fonttitle=\bfseries,title={#2\hfill #3},boxrule=0.8pt,arc=2mm,
  before skip=8pt,after skip=8pt}
\newtcolorbox{astuce}[1][Différenciation et besoins particuliers]{breakable,
  colback=yellow!12,colframe=yellow!50!black,fonttitle=\bfseries\small,
  title={#1},boxrule=0.5pt,arc=1mm,before skip=6pt,after skip=6pt}
\newtcolorbox{code}[1][Code]{breakable,colback=black!4,colframe=black!60,
  fonttitle=\bfseries\small,title={#1},boxrule=0.5pt,arc=1mm,fontupper=\ttfamily\small,
  before skip=6pt,after skip=6pt}
% Blocs de code coloré (écrire le code sans accents : comme dans l'éditeur)
\newtcblisting{arduino}[1][Code Arduino]{listing only,breakable,colback=black!4,
  colframe=black!60,fonttitle=\bfseries\small,title={#1},boxrule=0.5pt,arc=1mm,
  before skip=6pt,after skip=6pt,
  listing options={language=C++,basicstyle=\ttfamily\small,columns=fullflexible,
  keepspaces=true,showstringspaces=false,commentstyle=\color{green!40!black}\itshape,
  keywordstyle=\color{blue!70!black}\bfseries}}
\newtcblisting{python}[1][Code Python]{listing only,breakable,colback=black!4,
  colframe=black!60,fonttitle=\bfseries\small,title={#1},boxrule=0.5pt,arc=1mm,
  before skip=6pt,after skip=6pt,
  listing options={language=Python,basicstyle=\ttfamily\small,columns=fullflexible,
  keepspaces=true,showstringspaces=false,commentstyle=\color{green!40!black}\itshape,
  keywordstyle=\color{blue!70!black}\bfseries}}
\newcommand{\entete}[2]{%
  \begin{tcolorbox}[colback=#1!10,colframe=#1!70!black,boxrule=0.8pt,arc=2mm]
  {\Large\bfseries #2}\hfill Nom : \rule{4cm}{0.3pt}\quad Date : \rule{2.2cm}{0.3pt}
  \end{tcolorbox}}

% Bandeau de titre : #1 titre, #2 semaine
\newcommand{\titreblock}[2]{%
\begin{tcolorbox}[colback=teal!10,colframe=teal!70!black,arc=2mm,boxrule=1pt]
\centering
{\LARGE\bfseries #1}\\[2mm]
{\large Robotique -- \gradelabel{} -- Plan de cours de 2 heures}\\[1mm]
Enseignant : Francis Ki \quad|\quad #2
\end{tcolorbox}}

% Tableau du minutage : 6 contenus (accroche, mini-leçon, guidée, autonome, mise en commun, exit ticket)
\newcommand{\minutage}[6]{%
\begin{tabularx}{\linewidth}{>{\bfseries}p{2.2cm}p{2.6cm}Y p{3cm}}
\toprule
Minutes & Étape & Contenu & Organisation\\
\midrule
0 -- 10 & Accroche & #1 & Classe entière\\
10 -- 30 & Mini-leçon & #2 & Classe entière\\
30 -- 60 & Pratique guidée & #3 & Collectif puis binômes\\
60 -- 95 & Pratique autonome & #4 & Individuel sur poste\\
95 -- 110 & Mise en commun & #5 & Classe entière\\
110 -- 120 & Exit ticket & #6 & Individuel sur poste\\
\bottomrule
\end{tabularx}}

% Étayage progressif : #1 aide 1, #2 aide 2, #3 aide 3
\newcommand{\aides}[3]{%
\textbf{Étayage progressif de l'aide.}
\begin{center}
\begin{tabularx}{\linewidth}{>{\bfseries}p{2.4cm}Y p{4.4cm}}
\toprule
Niveau d'aide & Contenu & Qui intervient\\
\midrule
Aide 1 & #1 & Le binôme seul\\
Aide 2 & #2 & Le binôme avec la carte d'aide\\
Aide 3 & #3 & L'enseignant\\
\bottomrule
\end{tabularx}
\end{center}}

% Les trois groupes de la pratique autonome : #1 tâche A, #2 tâche B, #3 tâche C
\newcommand{\groupes}[3]{%
\textbf{Consigne commune\,:} chaque élève ouvre le fichier de son groupe, le renomme avec son prénom et travaille seul sur son poste. L'enseignant passe d'abord auprès du groupe A pendant 10 min.

\medskip
\begin{tabularx}{\linewidth}{>{\bfseries}p{3cm}Y p{2.2cm}}
\toprule
Groupe & Tâche sur poste & Annexe\\
\midrule
A : en dessous du niveau & #1 & Annexe C\\
B : au niveau & #2 & Annexe D\\
C : au-dessus du niveau & #3 & Annexe E\\
\bottomrule
\end{tabularx}}

% Mise en commun : #1 groupe A, #2 groupe B, #3 groupe C, #4 conclusion
\newcommand{\miseencommun}[4]{%
\begin{etape}{teal}{5. Mise en commun}{95 -- 110 min}
\begin{enumerate}
 \item (2 min) Retour au calme, écrans mis en veille.
 \item (9 min) Trois passages courts de 3 min, avec valorisation :
 \begin{itemize}
  \item \textbf{Groupe A} : #1
  \item \textbf{Groupe B} : #2
  \item \textbf{Groupe C} : #3
 \end{itemize}
 \item (4 min) Conclusion : #4
\end{enumerate}
\end{etape}}

% Clôture : #1 question 1, #2 question 2
\newcommand{\cloture}[2]{%
\begin{etape}{gray}{6. Clôture : exit ticket}{110 -- 120 min}
Formulaire en ligne ou papier (annexe G) :
\begin{enumerate}
 \item #1
 \item #2
 \item Auto-évaluation : \emph{Je maîtrise / Je progresse / J'ai besoin d'aide.}
\end{enumerate}
Lire les réponses avant la séance suivante pour ajuster les groupes de besoin.
\end{etape}}

% Grille d'observation : #1 #2 #3 intitulés de colonnes
\newcommand{\grille}[3]{%
\subsection*{Grille d'observation (enseignant)}
\begin{tabularx}{\linewidth}{Y p{2.4cm} p{2.4cm} p{2.4cm} Y}
\toprule
Élève & #1 & #2 & #3 & Groupe suivant\\
\midrule
 & & & & \\ & & & & \\ & & & & \\ & & & & \\ & & & & \\
\bottomrule
\end{tabularx}}

% Signature de l'exit ticket
\newcommand{\autoeval}{Je me sens : \cb Je maîtrise \quad \cb Je progresse \quad \cb J'ai besoin d'aide}


\begin{document}
\titreblock{Conditions (Python)}{Semaine du 9 au 13 novembre 2026}

%==================================================================
\section{Informations générales}
\begin{tabularx}{\linewidth}{>{\bfseries}p{3.6cm}Y}
\toprule
Titre & Conditions : \texttt{if}, \texttt{elif}, \texttt{else}\\
Classe / effectif & Grade 10, 11 à 20 élèves, niveaux hétérogènes, élèves à besoins particuliers\\
Durée & 2 heures (120 min)\\
Lieu & Salle informatique, un poste par élève\\
Matériel et logiciels &
\begin{itemize}
 \item Un poste par élève avec l'éditeur \textbf{Mu} en mode \emph{Python 3}.
 \item Vidéoprojecteur ou tableau interactif avec Mu projeté.
 \item Un dossier partagé contenant les fichiers de départ \texttt{.py} des trois niveaux.
 \item Inspiration : unités «~At a crossroads~» et «~More branches~» du cours TeachComputing (binômes, codage en direct, exemples commentés).
\end{itemize}\\
Pré-requis &
Séance précédente : variables, types \texttt{int}, \texttt{float}, \texttt{str}, saisie avec \texttt{input} et conversion.\\
\bottomrule
\end{tabularx}

\begin{astuce}[Préparation avant le cours (environ 40 min)]
\begin{enumerate}
 \item Préparer les fichiers de départ : \texttt{niveauA.py} (programme à compléter), \texttt{niveauB.py} (vide avec consignes), \texttt{niveauC.py} (pierre-feuille-ciseaux à terminer).
 \item Vérifier que l'indentation par quatre espaces fonctionne dans Mu (la touche Tab est bien configurée).
 \item Créer le formulaire de sortie (annexe G) ou imprimer l'annexe.
\end{enumerate}
\end{astuce}

%==================================================================
\section{Objectifs d'apprentissage}
\begin{tcolorbox}[colback=gray!8,colframe=gray!60!black,arc=1mm,boxrule=0.5pt]
\textbf{Objectifs du programme}
\begin{itemize}
 \item Utiliser les structures conditionnelles.
 \item Construire un programme décisionnel.
\end{itemize}
\textbf{À la fin de la séance, l'élève est capable\,:}
\begin{itemize}
 \item d'écrire une comparaison (\texttt{==}, \texttt{!=}, \texttt{<}, \texttt{>}, \texttt{<=}, \texttt{>=}) et de la combiner avec \texttt{and}, \texttt{or}, \texttt{not}\,;
 \item de construire une structure \texttt{if / elif / else} avec une indentation correcte\,;
 \item de tester toutes les branches d'un programme avec un jeu d'essais.
\end{itemize}
\textbf{Compétences numériques\,:} indenter son code, lire une erreur de syntaxe, tester toutes les branches.
\end{tcolorbox}

%==================================================================
\section{Déroulé minuté (120 min)}
\minutage{L'accès au parc : qui peut monter ?}{Comparaison, \texttt{if}, \texttt{elif}, \texttt{else}}{Un programme majeur ou mineur en collectif, puis un tarif de cinéma en binômes}{Trois niveaux de besoin}{Valoriser chaque groupe}{Trois questions sur poste}

%------------------------------------------------------------------
\begin{etape}{orange}{1. Accroche : l'accès au parc d'attractions}{0 -- 10 min}
\begin{enumerate}
 \item (5 min) Mise en situation : un agent contrôle l'accès à une attraction. \emph{Si la taille est d'au moins 1,20~m, alors tu montes ; sinon tu attends.} Plusieurs élèves jouent les visiteurs avec des tailles différentes.
 \item (5 min) Compliquer : \emph{et si tu as moins de 12 ans, il faut un adulte avec toi.} Question-problème : \textbf{«~Comment un programme décide-t-il quoi faire quand il y a plusieurs cas ?~»}
\end{enumerate}
\end{etape}

%------------------------------------------------------------------
\begin{etape}{blue}{2. Mini-leçon : décider avec if, elif, else}{10 -- 30 min}
\textbf{Support\,:} annexe A. Codage en direct. Quatre temps de 5 min.

\textbf{Temps 1 : comparer.} Une comparaison donne \texttt{True} ou \texttt{False}.
\begin{center}
\begin{tabularx}{\linewidth}{>{\ttfamily}p{2cm}Y Y}
\toprule
Opérateur & Signification & Exemple\\
\midrule
== & égal à & \texttt{age == 15}\\
!= & différent de & \texttt{age != 15}\\
< > & plus petit, plus grand & \texttt{taille > 1.2}\\
<= >= & inférieur ou égal, supérieur ou égal & \texttt{note >= 10}\\
\bottomrule
\end{tabularx}
\end{center}
Attention : \texttt{=} range une valeur, \texttt{==} compare deux valeurs.

\textbf{Temps 2 : if et else.}
\begin{python}[Deux chemins]
age = int(input("Ton age : "))
if age >= 18:
    print("Majeur")
else:
    print("Mineur")
\end{python}
L'indentation (quatre espaces) montre ce qui appartient au \texttt{if}.

\textbf{Temps 3 : elif pour plus de deux cas.}
\begin{python}[Trois chemins]
if note >= 16:
    print("Tres bien")
elif note >= 10:
    print("Passable")
else:
    print("Insuffisant")
\end{python}

\textbf{Temps 4 : and, or, not.} \texttt{if age >= 12 and taille >= 1.2:} (les deux conditions). \texttt{if jour == "samedi" or jour == "dimanche":} (l'une ou l'autre).

\begin{astuce}
Colorer le bloc indenté sur le tableau. Faire lire chaque condition à voix haute («~si l'âge est supérieur ou égal à 18~»). Fournir un modèle avec les deux-points et l'indentation déjà posés pour les élèves qui oublient.
\end{astuce}
\end{etape}

%------------------------------------------------------------------
\begin{etape}{green}{3. Pratique guidée}{30 -- 60 min}
\textbf{Fiche\,:} annexe B. Programmation en binômes (pilote et navigateur, rôles échangés à mi-temps).

\textbf{3a. Majeur ou mineur en collectif (12 min).} Écrire en codage en direct le programme complet, puis tester deux âges (17 et 18) pour vérifier la limite.

\textbf{3b. Tarif de cinéma en binômes (18 min).} Écrire un programme : moins de 12 ans, 3~\$ ; de 12 à 64 ans, 8~\$ ; 65 ans et plus, 5~\$. Tester les cas limites (11, 12, 64, 65).

\aides{Le programme du 3a et le tableau des opérateurs}{Questions guides : \emph{Quelle condition faut-il tester en premier ? Que se passe-t-il pour 12 ans exactement ?}}{L'enseignant tape la structure \texttt{if / elif / else} avec le binôme, qui complète}

\begin{astuce}
Faire dessiner sur papier les trois zones d'âge sur une droite avant de coder. Rappeler que Python lit les conditions dans l'ordre et s'arrête à la première qui est vraie.
\end{astuce}
\end{etape}

%------------------------------------------------------------------
\begin{etape}{violet}{4. Pratique autonome différenciée}{60 -- 95 min}
\groupes{Fichier \texttt{niveauA.py} à compléter : programme de mention à partir d'une note, avec les mots \texttt{if}, \texttt{elif}, \texttt{else}, \texttt{>=} fournis en banque. Exécuter avec trois notes (8, 12, 17) et noter l'affichage.}{Écrire un programme de mention complet (insuffisant, passable, assez bien, bien, très bien) puis un test d'année bissextile. Tester avec un jeu d'essais écrit à l'avance.}{Terminer le jeu pierre-feuille-ciseaux contre l'ordinateur (utiliser \texttt{import random} et \texttt{random.choice}), avec des conditions combinées par \texttt{and} et \texttt{or}. Vérifier les neuf combinaisons dans un tableau de tests.}

\medskip
\textbf{Plan B si un élève termine tôt\,:} ajouter un message de victoire personnalisé qui dépend du prénom saisi.
\end{etape}

%------------------------------------------------------------------
\miseencommun{un élève lit son programme et explique l'ordre des conditions. Valoriser l'explication.}{un binôme présente son jeu d'essais et la limite testée. Valoriser les cas limites.}{un élève joue contre son programme et explique la logique. Valoriser la créativité.}{\textit{un programme décisionnel teste des conditions dans l'ordre et exécute la première qui est vraie ; il faut tester toutes les branches.}}

\cloture{Quelle est la différence entre \texttt{=} et \texttt{==} ?}{Que s'affiche-t-il pour \texttt{note = 10} dans le programme de mention de l'annexe A ?}

%==================================================================
% ANNEXES
%==================================================================
\newpage
\section*{Annexe A : Support résumé de la mini-leçon}
\entete{blue}{Conditions en Python}

\begin{tcolorbox}[colback=yellow!10,colframe=yellow!50!black]
\texttt{if} teste une condition. \texttt{elif} ajoute un cas. \texttt{else} traite tous les autres cas. L'\textbf{indentation} (4 espaces) et les \textbf{deux-points} sont obligatoires.
\end{tcolorbox}

\begin{python}[Modèle]
note = int(input("Note : "))
if note >= 16:
    print("Tres bien")
elif note >= 10:
    print("Passable")
else:
    print("Insuffisant")
\end{python}

\begin{tabularx}{\linewidth}{>{\ttfamily}p{2.4cm}Y}
\toprule
Opérateur & Signification\\
\midrule
== != & égal, différent\\
< > <= >= & comparaisons\\
and & les deux conditions sont vraies\\
or & au moins une est vraie\\
not & inverse la condition\\
\bottomrule
\end{tabularx}

%==================================================================
\newpage
\section*{Annexe B : Fiche de pratique guidée}
\entete{green}{Pratique guidée : décider}

\subsection*{3a. Majeur ou mineur}
\begin{python}[À compléter pendant le codage en direct]
age = int(input("Ton age : "))
if age _____ 18:
    print("Majeur")
_____:
    print("Mineur")
\end{python}

\subsection*{3b. Tarif de cinéma}
\begin{tabularx}{\linewidth}{>{\bfseries}p{2.4cm}Y Y Y}
\toprule
Âge testé & Tarif attendu & Tarif obtenu & Conforme ?\\
\midrule
11 & & & \\ 12 & & & \\ 64 & & & \\ 65 & & & \\
\bottomrule
\end{tabularx}

\begin{tcolorbox}[colback=yellow!10,colframe=yellow!50!black,title={Cartes d'aide à utiliser dans l'ordre},fonttitle=\bfseries\small]
\textbf{Aide 1 :} regarde le programme du 3a et le tableau des opérateurs.\\
\textbf{Aide 2 :} Quelle condition faut-il tester en premier ? Que se passe-t-il pour 12 ans exactement ?\\
\textbf{Aide 3 :} demande à l'enseignant de taper la structure \texttt{if / elif / else} avec toi.
\end{tcolorbox}

%==================================================================
\newpage
\section*{Annexe C : Niveau A (en dessous du niveau)}
\entete{orange}{Niveau A : je complète la mention}

\textbf{Banque de mots :} if, elif, else, >=, 16, 10.

\begin{python}[niveauA.py]
note = int(input("Note : "))
_____ note >= 16:
    print("Tres bien")
_____ note _____ 10:
    print("Passable")
_____:
    print("Insuffisant")
\end{python}

\subsection*{Activité : trois essais}
\begin{tabularx}{\linewidth}{>{\bfseries}p{3cm}Y Y}
\toprule
Note saisie & Affichage prévu & Affichage obtenu\\
\midrule
8 & & \\ 12 & & \\ 17 & & \\
\bottomrule
\end{tabularx}

\subsection*{Activité : relie}
\texttt{==} \quad $\bullet$ \quad égal \qquad \texttt{!=} \quad $\bullet$ \quad différent \qquad \texttt{>=} \quad $\bullet$ \quad supérieur ou égal

%==================================================================
\newpage
\section*{Annexe D : Niveau B (au niveau)}
\entete{blue}{Niveau B : mention et année bissextile}

\subsection*{Mission 1 : mention}
Insuffisant (moins de 10), passable (10 à 11), assez bien (12 à 13), bien (14 à 15), très bien (16 et plus).

\subsection*{Mission 2 : année bissextile}
Une année est bissextile si elle est divisible par 4 mais pas par 100, ou si elle est divisible par 400. Indice : utilise \texttt{\%} et \texttt{and}, \texttt{or}.

\subsection*{Jeu d'essais (à écrire avant de coder)}
\begin{tabularx}{\linewidth}{>{\bfseries}p{3cm}Y Y Y}
\toprule
Entrée & Résultat attendu & Résultat obtenu & Conforme ?\\
\midrule
 & & & \\ & & & \\ & & & \\ & & & \\
\bottomrule
\end{tabularx}

%==================================================================
\newpage
\section*{Annexe E : Niveau C (au-dessus du niveau)}
\entete{violet}{Niveau C : pierre, feuille, ciseaux}

\begin{python}[niveauC.py : début du programme]
import random
choix = ["pierre", "feuille", "ciseaux"]
joueur = input("pierre, feuille ou ciseaux ? ")
ordi = random.choice(choix)
print("L'ordinateur a joue :", ordi)
# a toi de decider qui gagne
\end{python}

\textbf{Règles :} pierre bat ciseaux, ciseaux bat feuille, feuille bat pierre.

\begin{tabularx}{\linewidth}{>{\bfseries}p{2.4cm}Y Y Y}
\toprule
Joueur & Ordinateur & Résultat attendu & Conforme ?\\
\midrule
pierre & pierre & & \\ pierre & feuille & & \\ pierre & ciseaux & & \\
feuille & pierre & & \\ feuille & feuille & & \\ feuille & ciseaux & & \\
ciseaux & pierre & & \\ ciseaux & feuille & & \\ ciseaux & ciseaux & & \\
\bottomrule
\end{tabularx}

\textbf{Défi bonus :} refuse une saisie invalide en affichant un message.

%==================================================================
\newpage
\section*{Annexe F : Corrections et attendus (enseignant)}

\subsection*{Annexe B}
3a : \texttt{>=} et \texttt{else}. 3b : \texttt{if age < 12: 3 \$ ; elif age < 65: 8 \$ ; else: 5 \$}. Résultats : 11 donne 3 ; 12 donne 8 ; 64 donne 8 ; 65 donne 5.

\subsection*{Annexe C}
\texttt{if}, \texttt{elif}, \texttt{>=}, \texttt{else}. 8 : Insuffisant ; 12 : Passable ; 17 : Tres bien.

\subsection*{Annexe D}
Mention : \texttt{if / elif} dans l'ordre décroissant (16, 14, 12, 10). Bissextile : \texttt{(annee \% 4 == 0 and annee \% 100 != 0) or annee \% 400 == 0}. Essais : 2024 oui ; 1900 non ; 2000 oui ; 2023 non.

\subsection*{Annexe E}
Programme : égalité si \texttt{joueur == ordi}, puis \texttt{(joueur == "pierre" and ordi == "ciseaux") or (joueur == "ciseaux" and ordi == "feuille") or (joueur == "feuille" and ordi == "pierre")} pour une victoire, sinon défaite. Trois égalités, trois victoires, trois défaites.\\
\textbf{Critères de réussite :} les neuf cas testés, parenthèses correctes.

%==================================================================
\newpage
\section*{Annexe G : Exit ticket et grille d'observation}
\entete{gray}{Exit ticket : conditions}

\begin{enumerate}
 \item Quelle est la différence entre \texttt{=} et \texttt{==} ?\\ \lig
 \item Que s'affiche-t-il pour \texttt{note = 10} dans le programme de mention de l'annexe A ?\\ \lig
 \item \autoeval
\end{enumerate}
\textbf{Attendus :} 1. \texttt{=} range une valeur, \texttt{==} compare. 2. \texttt{Passable}, car \texttt{10 >= 10}.

\bigskip
\grille{Indentation correcte}{Cas limites testés}{Programme complet}

\end{document}
